% ============================================================
%  MODULE VI — Trading Mechanics
% ============================================================

\chapter{Trading Mechanics}

% ────────────────────────────────────────────────────────────
\section{Trading Timeframes}

Different traders operate on completely different timescales with different objectives:

\begin{center}
\begin{tabular}{>{\bfseries}p{3.5cm} p{3.5cm} p{6cm}}
\toprule
Style & Hold Period & Approach \\
\midrule
Scalping & Seconds -- Minutes & Tiny margins, massive volume; low retail suitability \\
Day Trading & Intraday only & No overnight positions; requires intense focus \\
Swing Trading & Days -- Weeks & Catch medium moves; most accessible to part-time traders \\
Position Trading & Weeks -- Months & Trend-following; fewer decisions required \\
Investing & Months -- Years & Fundamentals-driven; buy and hold \\
\bottomrule
\end{tabular}
\end{center}

\begin{notebox}
Most beginners should start with \textbf{swing trading or position trading}. Day trading and
scalping require exceptional skill, fast infrastructure, and very low commissions to be
profitable. Studies consistently show $>$70\% of day traders lose money.
\end{notebox}

% ────────────────────────────────────────────────────────────
\section{The Trading Plan}

Every trade should follow a pre-defined plan. Never trade without one.

A complete trade plan defines:
\begin{enumerate}[label=\textbf{\arabic*.}]
  \item \textbf{Setup / Thesis} — why does this trade make sense? What is the edge?
  \item \textbf{Entry criteria} — exact conditions that trigger entry (price level, indicator, event)
  \item \textbf{Stop loss} — where are you wrong? Maximum loss you accept on this trade.
  \item \textbf{Target(s)} — where do you take profit? Minimum first target.
  \item \textbf{Position size} — how many shares/contracts, given your stop and risk per trade.
  \item \textbf{Trade management} — do you scale out? Trail a stop? Let it run?
  \item \textbf{Invalidation} — what price/event would invalidate the entire thesis?
\end{enumerate}

% ────────────────────────────────────────────────────────────
\section{Position Sizing}

This is how much capital you allocate to a single trade. Critical for survival.

\begin{formulabox}
\textbf{Risk-based position sizing:}
\[
  \text{Position Size} = \frac{\text{Account Risk (\$)}}{\text{Stop Distance per Share (\$)}}
\]

\textbf{Account Risk:}
\[
  \text{Account Risk} = \text{Account Size} \times \text{Risk \% per Trade}
\]
\end{formulabox}

\begin{examplebox}
Account: \$20{,}000. Risk 1\% per trade. Entry: \$50. Stop: \$47.

\[
  \text{Account Risk} = \$20{,}000 \times 0.01 = \$200
\]
\[
  \text{Stop Distance} = \$50 - \$47 = \$3
\]
\[
  \text{Position Size} = \frac{\$200}{\$3} = 66 \text{ shares}
\]
Maximum loss on the trade = \$200, regardless of position size.
\end{examplebox}

\textbf{Common risk per trade for beginners: 0.5\%--1\%}. Professional traders often risk
0.25\%--0.5\% per trade.

% ────────────────────────────────────────────────────────────
\section{Risk-Reward Ratio}

\begin{keyterm}{Risk-Reward Ratio (RRR)}
The ratio of potential profit to potential loss on a trade.
\[
  RRR = \frac{\text{Target Distance}}{\text{Stop Distance}}
\]
\end{keyterm}

\begin{examplebox}
Entry: \$50. Stop: \$47 (risk = \$3). Target: \$59 (reward = \$9). \\
RRR = 9/3 = \textbf{3:1}
\end{examplebox}

With a 3:1 RRR, you only need to be right \textbf{25\% of the time} to break even.
Most professional traders target a minimum of \textbf{2:1 or 3:1} RRR.

\begin{formulabox}
\textbf{Breakeven Win Rate:}
\[
  W_{\min} = \frac{1}{1 + RRR}
\]
At 2:1 RRR: $W_{\min} = 1/3 \approx 33\%$. At 3:1: $W_{\min} = 25\%$.
\end{formulabox}

% ────────────────────────────────────────────────────────────
\section{Expected Value (EV)}

The single most important concept for evaluating a trading strategy:

\begin{formulabox}
\[
  EV = (P_{\text{win}} \times \text{Avg Win}) - (P_{\text{loss}} \times \text{Avg Loss})
\]
\end{formulabox}

\begin{examplebox}
Win rate = 40\%. Average win = \$300. Average loss = \$100.
\[
  EV = (0.40 \times 300) - (0.60 \times 100) = 120 - 60 = +\$60 \text{ per trade}
\]
This is a positive expectancy system — profitable over many repetitions.
\end{examplebox}

% ────────────────────────────────────────────────────────────
\section{Long vs. Short Selling}

\subsection{Going Long}
Buy shares, hold, sell later at higher price. Maximum loss = 100\% of capital invested.
Unlimited upside.

\subsection{Short Selling}
\begin{enumerate}[label=\arabic*.]
  \item \textbf{Borrow} shares from broker (pays a borrow fee).
  \item \textbf{Sell} those shares at current market price.
  \item \textbf{Buy back} (``cover'') at a lower price later.
  \item \textbf{Return} shares to lender and keep the difference.
\end{enumerate}

\begin{warningbox}
Short selling has \textbf{theoretically unlimited risk} — a stock can rise to infinity.
Losses are not capped. Additionally, a \textbf{short squeeze} can occur when many short sellers
are forced to buy simultaneously, catapulting the price rapidly (e.g., GameStop 2021).
\end{warningbox}

% ────────────────────────────────────────────────────────────
\section{Leverage and Margin}

\subsection{How Margin Works}
Margin allows you to borrow from your broker to control a larger position.
A \textbf{2:1 margin} means you can control \$20{,}000 of assets with \$10{,}000 of capital.

\subsection{Margin Call}
If losses erode your equity below the broker's \textbf{maintenance margin requirement},
you receive a margin call: deposit more funds or positions are forcibly liquidated —
often at the worst possible time.

\begin{warningbox}
Leverage is not a tool for beginners. Master position sizing and risk management
with \emph{no leverage} first. Many traders blow up accounts by using leverage
before they have developed the discipline and edge to use it responsibly.
\end{warningbox}

% ────────────────────────────────────────────────────────────
\section{Trading Psychology}

Markets are 50\% mechanics and 50\% psychology. Common destructive patterns:

\begin{itemize}
  \item \textbf{Revenge trading} — taking impulsive trades to recover a loss immediately.
        Almost always leads to larger losses.
  \item \textbf{FOMO} (Fear of Missing Out) — chasing price after missing an entry.
  \item \textbf{Letting losers run} — not respecting stop losses, hoping it will come back.
  \item \textbf{Cutting winners short} — exiting too early out of fear of giving back gains.
  \item \textbf{Overtrading} — taking too many trades; dilutes edge and increases commission costs.
  \item \textbf{Confirmation bias} — seeking information that confirms your position,
        ignoring contradictory evidence.
\end{itemize}

\begin{notebox}
A losing trade \emph{executed according to plan} is a good trade. A winning trade taken
outside your rules is a bad trade. \textbf{Process over outcome.} Trading is a probability game
measured over hundreds of trades, not individual outcomes.
\end{notebox}
